% !TEX root = ../../main.tex

\section{研究意义与国内外研究现状}

\subsection{研究意义}

本文工作的研究意义主要体现在以下几个方面。

第一，本文面向听障辅助训练场景构建了一个可运行的移动端训练系统。
系统将课程学习、数字人动作示范、手势训练、识别反馈和学习记录等功能组织在统一的移动端应用中，使用户能够按照较清晰的流程完成训练任务。
相比单纯的视频教学或静态课程展示，该系统更强调训练过程中的交互性和反馈闭环，有助于提升用户训练过程的连续性和可理解性。

第二，本文将手势识别能力与实际应用系统进行了工程化结合。
手势识别研究如果仅停留在离线模型训练和数据集测试阶段，往往难以直接体现其应用价值。
本文通过 Python 识别服务提供实时识别、数据采集、模型训练、模型运行时管理和句子视频识别等能力，并通过 Spring Boot 服务端与 HarmonyOS 移动端进行联动，使识别模型能够参与真实业务流程，包括训练评分、识别结果展示和句子级辅助交流结果生成。

第三，本文引入了有限词表下的句子视频识别与语义修正流程。
完整开放域连续手语识别和手语翻译对数据规模、标注质量、语言建模和模型泛化能力要求较高，超出了普通毕业设计系统的合理实现范围。
本文没有将系统目标设定为开放域连续手语翻译，而是在固定词表和受限场景下探索“视频输入—动作窗口分析—词级候选序列—中文映射—语义修正—句子级展示”的技术链路。
该设计既避免了对完整手语翻译能力的过度承诺，又能够验证句子级识别结果在辅助交流场景中的可用性。

第四，本文构建了面向模型维护的后台管理能力。
系统通过 Vue 后台管理端和 Spring Boot 服务端提供手势分类、手势资源、样本数据、训练任务和模型版本管理等功能，使识别模型的训练、发布和运行状态能够被统一维护。
这种设计将普通用户训练流程与开发者或管理员的模型维护流程相对分离，有助于提高系统的可维护性，也为后续扩展更大词表、更高质量样本和更多模型版本提供了基础。

\subsection{听障辅助训练与手语数字人国内外研究现状}

随着无障碍信息服务和人工智能技术的发展，面向听障群体的辅助系统逐渐从单一的信息提示工具，发展为融合课程资源、手语展示、动作识别、语义转换和交互反馈的综合系统。
听障用户在日常生活中不仅需要获取外部语音信息，也需要通过手语、文字、图像和视频等多种形式完成表达与交流。
联合国关于国际手语日的说明指出，手语是具有完整结构的自然语言，不同于口语语言\cite{un-sign-languages-day}。
因此，如何利用数字化方式降低沟通成本、提升训练效率，是听障辅助系统研究中的重要问题。

在手语教学和无障碍表达方向，手语数字人是一类具有代表性的技术路线。
与真人录制视频相比，数字人动作资源具有可复用、可组合和可统一管理等优点，适合用于课程教学、公共服务信息展示和标准动作示范等场景。
早期手语数字人研究常依赖符号化手语描述体系，其中 HamNoSys 是一种用于描述手语动作参数的书写系统，可表示手形、位置、运动方向及非手动特征等内容\cite{hanke-hamnosys-2004}。
在此基础上，SiGML 作为面向虚拟人动画的手势标记语言，被用于驱动手语虚拟人完成动作展示\cite{elliott-sigml-2004}。
JASigning 等系统进一步将 SiGML 与虚拟人动画结合，使浏览器或客户端能够展示由符号描述驱动的手语动作\cite{uea-jasigning}。
这些研究为手语数字人的动作描述、资源组织和动画展示提供了重要基础。

近年来，手语数字人研究开始从规则驱动和符号驱动逐步扩展到数据驱动和智能生成方向。
Kipp 等人的研究指出，手语虚拟人的核心难点不仅在于手部动作是否正确，还包括面部表情、身体姿态、节奏韵律以及听障用户对虚拟人表达效果的可理解性评价\cite{kipp-sign-avatar-2011}。
Bragg 等人也强调，手语识别、生成和翻译并不是单纯的计算机视觉问题，而是涉及语言学、计算机图形学、人机交互和听障文化理解的综合问题\cite{bragg-sign-language-ai-2019}。
SignON 项目则进一步探索了移动端手语识别、语音或文本处理和三维虚拟人表达相结合的系统框架，体现出手语辅助交流系统向多模态、移动化和用户参与式设计发展的趋势\cite{signon-project}。

从现有研究可以看出，手语数字人具有较好的教学展示和辅助交流潜力，但其工程化落地仍面临一定难点。
一方面，若采用符号化手语描述，需要构建规范的手语资源库，并处理不同手语体系、动作变体和非手动特征表达问题；
另一方面，若采用深度学习生成方式，则需要大规模高质量手语数据，并解决动作自然性、时序连续性和用户可理解性评价等问题。
对于毕业设计阶段的应用系统而言，直接实现完整的文本到手语自动生成或连续手语翻译难度较高，而将数字人作为标准动作展示、训练示范和课程资源载体，则更符合实际开发条件。

基于上述研究现状，本文所设计的 HearBridge 系统没有将重点放在完整的手语自动翻译生成上，而是将手语数字人作为移动端训练场景中的动作展示资源使用。
系统通过后台维护手势资源、课程内容和数字人资源地址，移动端在训练前向用户展示对应动作示范，使用户能够在较直观的视觉引导下进行手势学习与练习。
同时，系统结合实时手势识别服务，对用户练习动作进行识别和反馈，从而形成“标准动作展示—用户模仿训练—识别结果反馈”的训练流程。
这种设计避开了完整手语生成系统对大规模语料和复杂动画合成的高度依赖，更适合在实际移动应用中进行实现与验证。

\subsection{手势识别与手语识别相关研究现状}

手势识别与手语识别是听障辅助训练系统中的关键技术。
按照任务复杂度划分，相关任务通常可以分为孤立手势识别、连续手语识别和手语翻译。
孤立手势识别主要针对单个动作或单个词汇进行分类，适合用于课程训练、动作考核和基础指令识别等场景；
连续手语识别需要从连续视频中识别多个手语词汇序列，面临动作边界不明确、过渡动作复杂、语速差异和上下文依赖等问题；
手语翻译则进一步要求将视觉动作序列转换为自然语言表达，需要同时处理视觉识别、手语语法和自然语言生成问题\cite{bragg-sign-language-ai-2019}。
因此，从工程实现难度和毕业设计目标看，孤立手势识别适合作为系统基础能力，有限词表句子视频识别可以作为连续手语识别方向的受限探索，而开放域连续手语翻译更适合作为后续扩展目标。

在技术方法上，早期手势识别研究较多依赖人工特征和传统分类器。
随着深度学习的发展，卷积神经网络、循环神经网络、图卷积网络和 Transformer 等方法逐渐被用于手势和手语识别任务。
Koller 等人提出的 CNN-HMM 混合方法将卷积神经网络的视觉特征提取能力与隐马尔可夫模型的时序建模能力结合，用于连续手语识别任务\cite{koller-deep-sign-2018}。
Camgoz 等人进一步将神经机器翻译思想引入手语翻译任务，强调手语并不是简单的手势标签序列，而具有独立语法和语言结构\cite{camgoz-neural-slt-2018}。
在后续研究中，Sign Language Transformers 通过 Transformer 架构联合建模连续手语识别和手语翻译任务，体现出手语识别研究从单纯动作分类向视觉语言理解方向发展的趋势\cite{camgoz-slt-transformers-2020}。

直接基于原始视频进行识别能够保留较丰富的视觉信息，但也更容易受到背景、光照、摄像头角度、人体位置和设备性能等因素影响。
对于移动端训练系统而言，如果直接使用复杂视频模型，系统部署和实时运行成本较高，也不利于在较短开发周期内完成稳定联调。
基于关键点的识别方案为轻量化手势识别提供了可行路径。
MediaPipe 等视觉感知工具能够从图像帧中提取手部和人体姿态关键点，将原始视频转换为结构化坐标序列\cite{lugaresi-mediapipe-2019}。
其中，MediaPipe Hands 能够在设备端进行实时手部关键点跟踪，适合实时交互和移动端相关场景\cite{zhang-mediapipe-hands-2020}。
与直接处理整帧图像相比，关键点特征能够减少背景信息干扰，降低输入维度，并便于构造手型、相对位置、关节角度和动态变化等时序特征。

数据集建设也是手语识别研究中的重要基础。
WLASL 数据集提供了较大规模的词级美国手语视频数据，为孤立词级手语识别研究提供了重要参考\cite{li-wlasl-2020}。
在中文手语方向，CSL-Daily 数据集面向日常生活场景，提供了连续手语翻译所需的手语视频、gloss 标注和中文文本翻译\cite{zhou-csl-daily-2021}。
近期 NationalCSL-DP 数据集进一步围绕中国国家通用手语构建大规模双视角孤立手语识别数据集，包含较多手语词汇和双视角视频数据，为中文孤立手语识别研究提供了更丰富的数据基础\cite{jin-national-csl-dp-2025}。
这些数据集说明，手语识别研究正在逐步从小规模、实验室环境下的动作分类，发展到更大词表、更复杂场景和更强泛化能力的模型训练。

不过，现有研究成果在移动端听障训练系统中的直接应用仍存在一定限制。
首先，大规模公开数据集往往面向特定国家或地区的手语体系，不同手语之间不能简单通用；
其次，连续手语识别和手语翻译对数据量、标注质量和模型规模要求较高，直接部署到普通移动训练应用中存在较大工程压力；
再次，真实用户训练场景中的光照、摄像头角度、身体姿态和动作规范程度差异较大，容易影响识别稳定性。
因此，对于面向课程训练的应用系统而言，采用有限动作集合、固定训练项、关键点特征和轻量分类模型的方案更具有可实现性。

在本文实现中，系统首先完成固定训练项下的单词级手势识别能力，并将其用于移动端训练反馈；
随后进一步实现有限词表下的句子视频识别实验。
该实验以短视频片段为输入，通过动作窗口分析和滑动窗口推理获得词级候选序列，再结合中文映射和语义修正方法生成句子级结果。
语义修正部分利用大语言模型对候选词序列进行整理和表达修正，相关技术基础来自近年来大语言模型在自然语言理解与生成方面的能力提升\cite{deepseek-v3-2024}。
需要说明的是，该部分工作仍属于有限词表、受限场景下的句子级识别探索，不等同于完整开放域连续手语翻译；
但它验证了从视频识别结果到句子级展示的基本链路，为后续扩展连续手语识别和辅助交流功能提供了实验基础。
